\section{Introducción}

El presente trabajo práctico consiste en el diseño sistemático, 
simulación e implementación física de un filtro activo que 
cumpla con una plantilla de requerimientos específicos de 
atenuación. \\

Para este diseño, se utilizará la técnica de aproximación 
mediante polinomios de Chebychev, la cual permite obtener una 
transición más abrupta entre la banda de paso y la de rechazo a 
costa de permitir un rizado (ripple) controlado en la banda de 
paso. La metodología empleada abarca desde el cálculo analítico 
y la obtención de la función de transferencia mediante 
herramientas computacionales (Matlab), hasta la síntesis 
circuital utilizando topologías bicuadráticas de segundo orden, 
permitiendo así un análisis detallado de la sensibilidad y el 
comportamiento real frente a tolerancias de componentes.

\newpage
\section{Objetivos}

Los objetivos principales de esta actividad son: 

\begin{itemize}
    \item Sintetizar un circuito basado en amplificadores 
    operacionales que satisfaga la planilla de requerimientos 
    suministrada.
    \item Familiarizarse con el uso de herramientas de software 
    (Matlab/Pspice) para la aproximación de funciones de 
    atenuación y simulación de circuitos.
    \item Analizar el impacto de las tolerancias de los 
    componentes mediante simulaciones de Montecarlo y el cálculo 
    de sensibilidad de los parámetros del filtro ($\omega_p$ y 
    $Q_p$).
    \item Contrastar los resultados experimentales obtenidos en 
    laboratorio con las predicciones teóricas y las simulaciones 
    previas.
\end{itemize}